% ==========================================================================
%  PERGAMENA — template LaTeX per atti solenni in stile pergamena antica
%  Compilare con:  pdflatex pergamena.tex   (oppure latexmk -pdf)
%  Funziona anche su Overleaf (compilatore pdfLaTeX).
% ==========================================================================
\documentclass[11pt,a4paper]{article}

% --- Lingua e codifica -----------------------------------------------------
\usepackage[T1]{fontenc}
\usepackage[utf8]{inputenc}
\usepackage[italian]{babel}

% --- Caratteri -------------------------------------------------------------
\usepackage[oldstyle,proportional]{ebgaramond} % corpo del testo
\usepackage{cinzel}                            % titoli lapidari: \cinzel, \cinzel\itshape = Decorative
\usepackage{microtype}

% --- Impaginazione ---------------------------------------------------------
\usepackage[a4paper,top=3.6cm,bottom=3.8cm,left=3.4cm,right=3.4cm]{geometry}
\usepackage{xcolor}
\usepackage{tikz}
\usetikzlibrary{decorations.pathmorphing,calc,shadings}
\usepackage{pgfornament}
\usepackage{eso-pic}
\usepackage{lettrine}
\usepackage{enumitem}
\usepackage{fancyhdr}
\usepackage{setspace}
\usepackage{needspace}

% ==========================================================================
%  PALETTE
% ==========================================================================
\definecolor{pergChiara}{HTML}{F4E4BC}   % centro del foglio
\definecolor{pergMedia}{HTML}{E2C88E}    % tono intermedio
\definecolor{pergScura}{HTML}{B98B4E}    % bordi ingialliti
\definecolor{bruciato}{HTML}{3A2414}     % fondo dietro al foglio
\definecolor{macchia}{HTML}{8A5A2B}      % macchie d'umidità
\definecolor{inchiostro}{HTML}{3B2412}   % testo
\definecolor{inchiostroTenue}{HTML}{8C6A48}
\definecolor{rubrica}{HTML}{8E1B12}      % rosso dei capilettera (rubricatura)
\definecolor{ceralacca}{HTML}{9B1B1B}
\definecolor{ceralaccaScura}{HTML}{5E0D0D}

\color{inchiostro}

% ==========================================================================
%  SFONDO A PERGAMENA (disegnato su ogni pagina)
% ==========================================================================
\newcommand{\sfondopergamena}{%
\begin{tikzpicture}[remember picture,overlay]
  % semi diversi per ogni pagina, ma riproducibili
  \pgfmathsetseed{\number\numexpr 1789+97*\value{page}\relax}
  % 1) fondo scuro "bruciato" dietro al foglio
  \fill[bruciato] (current page.south west) rectangle (current page.north east);
  % 2) il foglio, con bordo irregolare e sfumatura radiale
  \def\bordo{([shift={(5mm,5mm)}]current page.south west) rectangle
             ([shift={(-5mm,-5mm)}]current page.north east)}
  \shade[inner color=pergChiara, outer color=pergScura,
         decorate, decoration={random steps,segment length=2.5mm,amplitude=1.1mm}]
         \bordo;
  % 3) vignettatura: margini più scuri e consumati
  \begin{scope}
    \clip[decorate, decoration={random steps,segment length=2.5mm,amplitude=1.1mm}] \bordo;
    \foreach \w/\o in {28mm/0.06, 18mm/0.07, 10mm/0.09, 5mm/0.12, 2mm/0.25}{
      \draw[macchia, line width=\w, opacity=\o,
            decorate, decoration={random steps,segment length=6mm,amplitude=2mm}] \bordo;
    }
    % 4) macchie d'umidità e d'età
    \foreach \i in {1,...,22}{
      \pgfmathsetmacro{\x}{rnd*\paperwidth}
      \pgfmathsetmacro{\y}{rnd*\paperheight}
      \pgfmathsetmacro{\r}{4+rnd*28}
      \pgfmathsetmacro{\op}{0.03+rnd*0.07}
      \fill[macchia, opacity=\op,
            decorate, decoration={random steps,segment length=2mm,amplitude=1.5mm}]
        ([shift={(\x pt,\y pt)}]current page.south west) circle (\r mm);
    }
    % 5) grana / fibre della pelle
    \foreach \i in {1,...,450}{
      \pgfmathsetmacro{\x}{rnd*\paperwidth}
      \pgfmathsetmacro{\y}{rnd*\paperheight}
      \pgfmathsetmacro{\a}{rnd*360}
      \pgfmathsetmacro{\l}{0.3+rnd*1.6}
      \draw[macchia, opacity=0.18, line width=0.15pt]
        ([shift={(\x pt,\y pt)}]current page.south west) -- ++(\a:\l mm);
    }
  \end{scope}
  % 6) cornice a doppio filetto con ornamenti agli angoli
  \def\m{1.7cm}
  \draw[inchiostro!80, line width=1.1pt]
    ([shift={(\m,\m)}]current page.south west) rectangle ([shift={(-\m,-\m)}]current page.north east);
  \draw[inchiostro!70, line width=0.4pt]
    ([shift={(\m+1.6mm,\m+1.6mm)}]current page.south west) rectangle ([shift={(-\m-1.6mm,-\m-1.6mm)}]current page.north east);
  \node[anchor=north west,inner sep=0] at ([shift={(\m+0.6mm,-\m-0.6mm)}]current page.north west)
    {\pgfornament[width=2.2cm,color=inchiostro]{61}};
  \node[anchor=north east,inner sep=0] at ([shift={(-\m-0.6mm,-\m-0.6mm)}]current page.north east)
    {\pgfornament[width=2.2cm,color=inchiostro,symmetry=v]{61}};
  \node[anchor=south west,inner sep=0] at ([shift={(\m+0.6mm,\m+0.6mm)}]current page.south west)
    {\pgfornament[width=2.2cm,color=inchiostro,symmetry=h]{61}};
  \node[anchor=south east,inner sep=0] at ([shift={(-\m-0.6mm,\m+0.6mm)}]current page.south east)
    {\pgfornament[width=2.2cm,color=inchiostro,symmetry=c]{61}};
\end{tikzpicture}}

\AddToShipoutPictureBG{\sfondopergamena}

% ==========================================================================
%  INTESTAZIONI E PIÈ DI PAGINA
% ==========================================================================
\pagestyle{fancy}
\fancyhf{}
\renewcommand{\headrulewidth}{0pt}
\fancyfoot[C]{\small\textsc{— foglio \thepage{} —}}
\fancypagestyle{plain}{\fancyhf{}\fancyfoot[C]{\small\textsc{— foglio \thepage{} —}}}

% ==========================================================================
%  COMANDI DI COMODO
% ==========================================================================
\setstretch{1.12}
\setlength{\parindent}{1.5em}
\setlength{\parskip}{0.35em}

% Campo da compilare: [ testo ] in inchiostro sbiadito
\newcommand{\campo}[1]{\textcolor{inchiostroTenue}{[\,\textit{#1}\,]}}

% Separatore ornamentale
\newcommand{\fregio}[1][88]{%
  \par\medskip
  {\centering\pgfornament[width=4.5cm,color=inchiostro]{#1}\par}
  \medskip}

% Capolettera rubricato (rosso)
\renewcommand{\LettrineFontHook}{\initials\color{rubrica}}
\renewcommand{\LettrineTextFont}{\scshape}
\setcounter{DefaultLines}{3}
\newcommand{\capolettera}[2]{\lettrine[lines=3,findent=0.25em,nindent=0pt]{#1}{#2}}

% Titolo solenne
\newcommand{\titoloatto}[2]{%
  \thispagestyle{plain}
  \begin{center}
    {\cinzel\itshape\fontsize{30}{36}\selectfont #1\par}
    \vspace{0.6em}
    {\pgfornament[width=7cm,color=inchiostro]{84}\par}
    \vspace{0.4em}
    {\cinzel\large\textls[150]{#2}\par}
  \end{center}
  \vspace{1em}}

% Rubrica di sezione (es. "Premesso che")
\newcommand{\rubrica}[1]{%
  \par\Needspace{4\baselineskip}\bigskip
  {\centering\cinzel\color{rubrica}\large #1\par}
  \smallskip}

% Articoli numerati in romano
\newcounter{articolo}
\newcommand{\articolo}[1]{%
  \stepcounter{articolo}%
  \par\Needspace{7\baselineskip}\bigskip
  {\centering
    {\cinzel\color{rubrica}Articolo \Roman{articolo}\par}
    \vspace{0.15em}
    {\scshape\large #1\par}
    \vspace{0.1em}
    {\pgfornament[width=1.6cm,color=inchiostro!80]{89}\par}}
  \smallskip\noindent\ignorespaces}

% Commi numerati
\newlist{commi}{enumerate}{1}
\setlist[commi]{label=\arabic*.,leftmargin=2em,itemsep=0.25em,topsep=0.3em}

% Elenco delle premesse
\newlist{premesse}{enumerate}{1}
\setlist[premesse]{label=\textit{\alph*)},leftmargin=2em,itemsep=0.25em,topsep=0.3em}

% Sigillo di ceralacca: \sigillo{iniziali}
\newcommand{\sigillo}[1]{%
\begin{tikzpicture}[baseline=(c.center)]
  \pgfmathsetseed{4242}
  % nastro
  \fill[rubrica!85!black] (-0.35,-0.3) -- (-0.85,-2.0) -- (-0.55,-1.8) -- (-0.4,-2.1) -- (0.05,-0.4) -- cycle;
  \fill[rubrica!70!black] (0.35,-0.3) -- (0.9,-1.95) -- (0.6,-1.8) -- (0.42,-2.05) -- (-0.05,-0.4) -- cycle;
  % colatura di cera
  \shade[inner color=ceralacca, outer color=ceralaccaScura,
         decorate, decoration={random steps,segment length=1.6mm,amplitude=0.9mm}]
    (0,0) circle (1.25);
  % impronta del timbro
  \draw[ceralaccaScura, line width=1.2pt] (0,0) circle (0.92);
  \draw[ceralaccaScura!80, line width=0.4pt] (0,0) circle (0.82);
  \fill[ceralacca!85!white, opacity=0.25] (-0.35,0.45) ellipse (0.35 and 0.18);
  \node (c) at (0,0) {\cinzel\itshape\color{ceralacca!55!pergChiara}\Large #1};
\end{tikzpicture}}

% Firma: \firma{qualifica}{nome}
\newcommand{\firma}[2]{%
  \begin{minipage}[t]{0.4\textwidth}\centering
    \textsc{#1}\par\vspace{1.6cm}
    \rule{\linewidth}{0.4pt}\par
    \small\textit{#2}
  \end{minipage}}

% ==========================================================================
%  DOCUMENTO
% ==========================================================================
\begin{document}

\titoloatto{Atto Solenne}{di \campo{oggetto dell'atto}}

\capolettera{L}{'anno duemilaventisei}, il giorno \campo{diciotto} del mese di
\campo{novembre}, in \campo{luogo}, presso \campo{indirizzo}, innanzi a \campo{nome del testimone
o del pubblico ufficiale} e alla presenza dei testimoni idonei, noti e richiesti,
sono personalmente comparsi e si sono costituiti i signori:

\begin{premesse}
  \item \campo{Nome e Cognome}, nato a \campo{luogo} il \campo{data}, residente in
        \campo{indirizzo}, codice fiscale \campo{codice}, di seguito per brevità
        denominato \guillemotleft\textsc{la Parte Prima}\guillemotright;
  \item \campo{Nome e Cognome}, nato a \campo{luogo} il \campo{data}, residente in
        \campo{indirizzo}, codice fiscale \campo{codice}, di seguito per brevità
        denominato \guillemotleft\textsc{la Parte Seconda}\guillemotright;
\end{premesse}

\noindent di seguito congiuntamente denominati \guillemotleft le Parti\guillemotright\
e singolarmente \guillemotleft la Parte\guillemotright.

\rubrica{Premesso che}

\begin{premesse}
  \item la Parte Prima è titolare di \campo{diritto, bene o facoltà}, libero da
        pesi, vincoli, oneri e trascrizioni pregiudizievoli di qualsivoglia natura;
  \item la Parte Seconda ha manifestato il proprio interesse a \campo{finalità},
        alle condizioni e nei termini di cui al presente atto;
  \item è intendimento delle Parti regolare per iscritto i reciproci diritti e
        obblighi, a ogni effetto di legge e di ragione.
\end{premesse}

\rubrica{Tutto ciò premesso}

\noindent che forma parte integrante e sostanziale del presente atto, le Parti,
come sopra costituite, convengono e stipulano quanto segue.

\fregio

\articolo{Delle premesse}
Le premesse e gli eventuali allegati costituiscono parte integrante e sostanziale
del presente atto e ne informano l'interpretazione.

\articolo{Dell'oggetto}
\begin{commi}
  \item Con il presente atto la Parte Prima \campo{cede / concede / si obbliga a}
        in favore della Parte Seconda, che accetta, \campo{descrizione dell'oggetto}.
  \item Quanto sopra si intende nello stato di fatto e di diritto in cui attualmente
        si trova, con ogni accessione, pertinenza e servitù attiva e passiva.
\end{commi}

\articolo{Degli obblighi delle Parti}
\begin{commi}
  \item La Parte Prima si obbliga a \campo{obblighi}, con la diligenza del buon padre
        di famiglia.
  \item La Parte Seconda si obbliga a corrispondere la somma di \campo{importo}
        (\campo{importo in lettere}), entro e non oltre il \campo{termine}.
\end{commi}

\articolo{Della durata e del recesso}
Il presente atto ha efficacia dalla data della sua sottoscrizione e sino al
\campo{data}. È facoltà di ciascuna Parte recedere con preavviso scritto di
\campo{giorni} giorni, da comunicarsi a mezzo raccomandata con avviso di
ricevimento ovvero posta elettronica certificata.

\articolo{Della clausola penale}
In caso di inadempimento, anche parziale, la Parte inadempiente sarà tenuta a
corrispondere all'altra, a titolo di penale, la somma di \campo{importo}, salvo
il risarcimento del maggior danno ai sensi dell'art.~1382 del Codice civile.

\articolo{Del foro competente}
Per ogni controversia che dovesse insorgere in ordine all'interpretazione,
esecuzione o risoluzione del presente atto, le Parti eleggono quale foro
esclusivamente competente quello di \campo{città}.

\articolo{Delle disposizioni finali}
\begin{commi}
  \item Ogni modifica al presente atto dovrà risultare da atto scritto, a pena
        di nullità.
  \item Per quanto non espressamente previsto si rinvia alle norme del Codice
        civile e delle leggi speciali in materia.
\end{commi}

\fregio[85]

\noindent\textit{Fatto, letto, confermato e sottoscritto in \campo{luogo},
addì \campo{data}, in due originali, uno per ciascuna Parte.}

\vspace{1.2cm}
\noindent
\firma{La Parte Prima}{\campo{Nome e Cognome}}\hfill
\raisebox{-1.3cm}{\sigillo{AS}}\hfill
\firma{La Parte Seconda}{\campo{Nome e Cognome}}

\end{document}
